\textbf{Čtyři termoplasty:} PE, PP, PVC, PS. \textbf{Tři reaktoplasty:} fenoplasty PF, epoxidové pryskyřice EP, melaminoformaldehydové pryskyřice MF.
\begin{itemize}
\item \textbf{Polyethylen (PE):} nepolární semikrystalický polymer, lehký, chemicky odolný a málo nasákavý. Vysokotlaký PE má nižší hustotu a měkne kolem $105\,^{\circ}\mathrm C$, nízkotlaký je hustší a měkne kolem $130\,^{\circ}\mathrm C$. Je výborný izolant: $\varepsilon_r\approx2{,}25$--$2{,}30$, $\tan\delta\approx(1$--$3)\cdot10^{-4}$, $\rho_V\sim10^{15}\,\Omega\,\mathrm m$. Používá se na kabelové izolace včetně vf techniky.
\item \textbf{Epoxidové pryskyřice (EP):} vytvrzením vzniká prostorová síť. Mají dobrou přilnavost ke kovům, keramice a sklu, malou nasákavost a dobré elektrické vlastnosti: $\varepsilon_r\approx4$, $\tan\delta\sim10^{-3}$ při 100 Hz. Přednáška uvádí dlouhodobou tepelnou odolnost asi $160\,^{\circ}\mathrm C$. Použití: zalévací hmoty, lepidla a lamináty FR-4.
\end{itemize}
